% META: {"id": "TCOMPL-AIR-CO-009", "chapitre": "TCOMPL-CALCULS-AIRES", "type_objet": "corrige", "exercice_ref": "TCOMPL-AIR-EX-009", "capacites_codes": ["C5"], "sources_inspiration": [], "status": "generated"}
% BEGIN-VERIFY
% from sympy import symbols, diff, simplify, Function, integrate
% x = symbols('x')
% F, G = x ** 2, x ** 2 + 7
% assert simplify(diff(F, x) - 2 * x) == 0
% assert simplify(diff(G, x) - 2 * x) == 0
% assert simplify(diff(F - G, x)) == 0
% assert simplify(F - G) == -7
% for c in (-3, 0, 5):
%     assert simplify(diff(x ** 2 + c, x) - 2 * x) == 0
% END-VERIFY
\begin{corrige}{TCOMPL-AIR-EX-009}
\textbf{1.} $(F-G)'=F'-G'=f-f=0$ sur $I$.

\textbf{2.} Une fonction dérivable de dérivée nulle sur un \emph{intervalle}
est constante : il existe $c$ tel que $F-G=c$, c'est-à-dire $F=G+c$.
L'hypothèse d'intervalle est essentielle : sur une réunion de deux
intervalles disjoints, la constante pourrait différer sur chacun.

\textbf{3.} Pour $f(x)=2x$, $F(x)=x^2$ et $G(x)=x^2+7$ sont deux primitives,
et $F-G=-7$ est bien constante.
\end{corrige}
